\documentclass[11pt]{article}
\usepackage[margin=2.2cm]{geometry}
\usepackage{amsmath,amssymb}
\usepackage{graphicx}
\usepackage{booktabs}
\usepackage{longtable}
\usepackage{xcolor}
\usepackage{hyperref}
\usepackage{caption}
\usepackage{float}
\usepackage[htt]{hyphenat}
\usepackage{enumitem}

\emergencystretch=3em
\sloppy
\setlength{\tabcolsep}{4.5pt}

\newcommand{\Q}{\mathbb Q}
\newcommand{\Pp}{\mathbb P}
\newcommand{\E}{\mathbb E}
\newcommand{\R}{\mathbb R}

\title{A CRC-inspired, dynamically-consistent extension of the\\
DYSANOS FLEX statistical simulator: audit, tests, and a\\
working consistent-recalibration sequence}
\author{\texttt{quant\_data-main/quantdata/dysanos/flex/crc\_extension}}
\date{\today}

\begin{document}
\maketitle

\begin{abstract}
We audit the existing \texttt{quantdata.dysanos.flex} statistical option-market
simulator, formalise and test its implicit discrete-time pricing identity, run a
drift-feasibility (rank/conditioning) analysis for the simplest admissible-$\Q$
family, and implement the one component the audit found genuinely missing: a
consistent-recalibration (CRC) sequence with an exact, feasibility-checked,
price-preserving reset. Every claim is checked against the running code in a
freshly built, fully-dependency-resolved Python~3.12 environment (18/18 of the
existing repository's own FLEX tests pass unmodified), not asserted from the
accompanying DYSANOS~II manuscript alone. No real historical option book was
available in this environment (the local test fixtures cover only a handful of
scattered trading days); every reported number is therefore a machinery/structural
validation on synthetic-but-structurally-faithful data, not a claim about real
market dynamics -- stated explicitly at every reported number, not left implicit.
\end{abstract}

\tableofcontents

\section{Executive summary}
\label{sec:exec}

\begin{enumerate}
\item \textbf{Environment}: no working Python environment or trained checkpoint
  existed in this workspace beforehand. A Python~3.12 virtual environment was
  built from the repository's own \texttt{pyproject.toml}; the existing FLEX test
  suite (\texttt{test\_dysanos\_torch\_simulator.py}, \texttt{test\_dysanos\_market.py})
  passes 18/18 against local fixture data. This is the untouched reproducible
  baseline every subsequent number is compared against.
\item \textbf{Audit}: the target implementation is \texttt{quantdata/dysanos/flex/}
  (Sec.~\ref{sec:audit}). Its statistical transition (``the existing AR(1) law'')
  is a PCA-reduced, hierarchical, per-expiry-block VAR(1)
  (\texttt{weekly\_dynamics.py}), with \textbf{no} $\Q$-restriction anywhere in it.
  A pathwise discrete HJM/roll residual (\texttt{hjm\_roll\_residuals}) and a
  gradient-trained ``target-measure'' kernel already exist in the codebase --
  substantial prior work directly on point, used here rather than reinvented.
  The clearest gap: no price-preserving reset (CRC step E) exists anywhere in
  \texttt{flex/}.
\item \textbf{Pricing identity, tested (Sec.~\ref{sec:sec23})}: the raw fitted P
  law decisively fails the discrete martingale identity (t-statistic $16.9$ at
  100{,}000 paths, well-powered, not assumed). The existing trained
  target-measure kernel reduces but does not eliminate the residual (t-statistic
  still $8.9$ after training convergence).
\item \textbf{Drift feasibility, ranked (Sec.~\ref{sec:sec23})}: a linearised,
  additive mean-shift drift family is rank-deficient ($6$ of $8$ directions
  active) and, even within its own reachable subspace, reduces the targeted
  residual by only $69\%$, not $100\%$ -- reported precisely, not rounded to
  ``works'' or ``fails.''
\item \textbf{The CRC sequence, implemented (Sec.~\ref{sec:sec5})}: steps A--G are
  implemented exactly as specified. A structural fact is exploited and verified,
  not assumed: FLEX's own bounded-decoder parametrisation makes static
  admissibility automatic for \emph{any} raw state, so the constrained inverse
  (step E) reduces to an unconstrained optimisation that always decodes
  admissibly. The sequence converges cleanly across 4 resets on synthetic data;
  a deliberately implausible target is correctly reported as infeasible,
  confirming the feasibility check is not vacuous.
\item \textbf{What is not established here}: full $\Q$-admissibility for a richer
  family (Sec.~\ref{sec:deferred}), genuine historical P/Q compatibility fitting,
  and validation against real market data all require the production
  \texttt{QUANTDATA\_ROOT}/WRDS-backed environment, not available here. These are
  scoped precisely, not silently omitted.
\end{enumerate}

\section{Audit of the existing implementation}
\label{sec:audit}

\subsection{Target identified, and why}

\texttt{quantdata/dysanos/} contains several largely independent DYSANOS
subsystems: \texttt{flex/}, \texttt{grid/gendysanos}, \texttt{grid/rndysanos},
\texttt{implied/}, \texttt{localvolbysanos/}, \texttt{tangent\_levy/}. This
report targets \texttt{flex/}, the only subsystem whose language (``AR(1)'',
weekly, joint DLV+grid+smoothing+dynamics estimation) matches the assignment's
description exactly, and the only one with an explicit, already-present
pathwise HJM/roll-residual diagnostic directly on point for the pricing
identity (Sec.~\ref{sec:sec23}). No trained checkpoint exists in this
repository checkout; numbers below fit fresh against local fixture/synthetic
data, stated as such throughout.

\subsection{Latent state and decoder}

\texttt{bound\_native\_raw} maps an unconstrained raw torch tensor to
\texttt{NativeParameters(dlvsqs, forward\_vols, eta)}: \texttt{dlvsqs}
$(T{\times}N)$ bounded to $[0.01^2,3.0^2]$, \texttt{forward\_vols} $(T,)$
bounded to $[0.01,1.0]$, \texttt{eta} (fixed or bounded $[0.01,0.50]$) via
sigmoid/interval transforms. \texttt{FlexSanosDecoder} then runs the
Wissel/Magnus finite-grid recursion directly on \texttt{dlvsqs} (confirmed
from the weight construction \texttt{w\_plus}/\texttt{w\_minus}, not assumed
from naming).

\subsection{Two term structures, and the variance/std convention}

\textbf{Grid geometry} (\texttt{num\_inner\_strikes}, \texttt{nms\_range},
\texttt{num\_model\_expiries}, the calendar roll schedule) is fixed per run.
\textbf{Smoothing} (\texttt{eta}, \texttt{forward\_vols}) is, materially,
\emph{part of the fitted state} -- \texttt{forward\_vols} is jointly
estimated, not a fixed side channel. Traced directly through
\texttt{price\_first\_expiry\_calls\_from\_parameters} $\to$
\texttt{\_black\_call}: the kernel dispersion is
\texttt{forward\_vol*eta*sqrt(dt)}, passed to the argument literally named
\texttt{sqrt\_variance} and used directly (not squared) in $d_1$ -- a standard
deviation, with \texttt{eta} itself a \emph{dimensionless multiplier} on
\texttt{forward\_vol}, not a standard deviation on its own. This differs
materially from other SANOS conventions in this research programme (where
$\eta$ \emph{is} the smoothing standard deviation directly) and must be
preserved exactly when porting any restriction between them.

\subsection{Statistical transition, front-expiry law, roll conventions}

\texttt{weekly\_dynamics.py}: a PCA-reduced, per-expiry-block, hierarchical
VAR(1) (\texttt{HierarchicalBlockFit(intercept, own\_lag, previous\_lag)}) at
a fixed 5-business-day stride, ridge-regularised, with a joint innovation
covariance and a spot process that loads on the surface innovation via
\texttt{weekly\_spot\_front\_beta} plus an orthogonal residual. A
continuous-time generator $K=-\log(A)/\Delta t$ is computed and \emph{gated}
(condition-number ceiling, exponential round-trip check) -- a P-side
stability diagnostic only; $\Q$ does not appear anywhere in this file.
\texttt{fixed\_expiry\_ar1.py}/\texttt{last\_expiry\_ar1.py} handle the DLV
tower's calendar-exact roll; the front-expiry settlement payoff used in the
codebase's own audit is the ordinary intrinsic value
(\texttt{relu(R-k)}), not a grid-pinned substitute.

\subsection{Existing pricing-drift infrastructure (the central prior-work finding)}

\texttt{WeeklySanosSimulator.hjm\_roll\_residuals} already computes, pathwise,
exactly the discrete martingale/roll residual Sec.~\ref{sec:sec23} formalises
and tests: \texttt{R*c\_next[m-1](k/R) - c\_current[m](k)} for non-front
expiries, with the correct terminal payoff at the front. Its only caller,
\texttt{weekly\_measure\_training.py}, already implements a sequential-Monte-
Carlo/gradient-based fit of a ``target-measure'' kernel
(\texttt{NormalWeeklyMeasureKernel}/\texttt{GenericWeeklyMeasureKernel}) via
\texttt{hjm\_roll\_moment\_loss} with instrumented conditioning -- a
substantial prior attempt at $\Pp\to\Q$ construction, used directly in
Sec.~\ref{sec:sec23} rather than reinvented. A family of
\texttt{run\_*\_q\_existence.py} scripts investigates the same admissibility
question via Wasserstein-2 roll conditions against real historical fits; their
result artefacts live in a training cache not present in this checkout and
were not read here -- flagged explicitly rather than silently worked around.

\textbf{The clearest gap}: no price-preserving reset (an exact constrained
inverse under admissibility constraints, CRC step E) exists anywhere in
\texttt{flex/}. This is what Sec.~\ref{sec:sec5} builds.

\subsection{Reproducible baseline}

Python~3.12, fresh \texttt{.venv312} built from \texttt{pyproject.toml},
\texttt{QUANTDATA\_TEST\_MODE=1}: \texttt{pytest
tests/test\_dysanos\_torch\_simulator.py tests/test\_dysanos\_market.py} --
\textbf{18/18 pass}. No file outside \texttt{flex/crc\_extension/} was
modified.

\section{The discrete-time pricing identity: derived, tested, and ranked for feasibility}
\label{sec:sec23}

\subsection{The identity}

For the normalised decoder, $H(t,S_t,h_t;T,K)=S_t\cdot c(h_t;T-t,K/S_t)$.
Dividing $\E_\Q[H(t+\Delta,Z_{\text{next}};T,K)\mid Z_t]=H(t,Z_t;T,K)$ by $S_t$
and writing $R=S_{t+\Delta}/S_t$ gives, in normalised units,
\[
\E_\Q\bigl[R\cdot c(h_{\text{next}};T-t-\Delta,K/(S_tR))\bigr]=c(h_t;T-t,K/S_t).
\]
This is \emph{exactly} what \texttt{hjm\_roll\_residuals} computes pathwise --
confirmed by direct derivation from the identity above, not by trusting the
existing docstring: for a non-front model-expiry slab $m\geq2$ the tower rolls
one week, so the $(t,K)$ contract at index $m$ becomes the $(t+\Delta,K)$
contract at index $m-1$ next week, giving
\texttt{R*c\_next[m-1](k/R) - c\_current[m](k)}; at the front expiry
($m=1$) the correct terminal payoff \texttt{relu(R-k)} replaces
\texttt{c\_next[0]}. FLEX is a discrete-time simulator by construction, so
this discrete identity \emph{is} the operative pricing restriction, not a
discretisation of a separate continuous-time one.

\textbf{Scope}: only a declared finite protected book (3 model expiries
$\times$ 5 strikes $=15$ contracts) is tested below -- concept (a) of the
assignment. No claim is made, or should be inferred, about the entire smooth
surface (concept (b)).

\subsection{Tested, well-powered, on synthetic-but-structurally-valid data}

\begin{center}
\begin{tabular}{p{0.42\textwidth}rl}
\toprule
Law & $\max|t\text{-stat}|$ & mean $|$residual$|$ by expiry \\
\midrule
Raw fitted P (AR(1), untouched) & \textbf{16.9} & $[0.00056,\,0.0052,\,0.0067]$ \\
Trained target-measure kernel (600 ep., converged) & \textbf{8.9} & $[0.00045,\,0.0043,\,0.0053]$ \\
\bottomrule
\end{tabular}
\par\vspace{2pt}\noindent\footnotesize 15 contracts (3 expiries $\times$ 5 strikes), 100{,}000 paths, one weekly step.
\end{center}

The raw P law decisively fails martingality -- consistent with, and now
numerically verified rather than merely asserted by, the codebase's own
\texttt{hjm\_roll\_moment\_loss} docstring (``Phase 3 does not claim that its
historical identity measure already satisfies risk-neutral HJM
conditions''). The existing trained target-measure kernel, trained to loss
convergence, \emph{reduces but does not eliminate} the residual.

\subsection{Drift feasibility: rank, conditioning, residual}

Linearised discrete-time analogue of the continuous-time tangent-space
feasibility test used throughout this research programme: does an additive
mean-shift $\mu\in\R^{8}$ (state dimension) to the P-law's next-period state
cancel the mean roll residual across all 15 tested contracts simultaneously,
to first order? Jacobian computed via autograd at $\mu=0$ (60{,}000-path
Monte Carlo mean):
\[
\text{residual\_dim}=15,\quad \text{mu\_dim}=8,\quad
\boxed{\text{rank}=6},\quad \text{cond}=\infty,\quad
\boxed{\text{reduction fraction}=0.691}.
\]
The Jacobian is \textbf{rank-deficient} (the longest-dated expiry block's 2
coordinates have exactly zero first-order effect on the tested residual
within one week -- confirmed directly: the corresponding two components of
the least-squares solution are exactly $0$, not merely small), and even
restricted to its own 6-dimensional reachable subspace, an additive mean-shift
cannot fully cancel a 15-dimensional residual target ($69\%$ reduction, not
$100\%$).

\textbf{This is not evidence that no richer family can work} -- a full
covariance tilt, a larger instrument set, or a structural product
construction (Route~B in the sense of this research programme's other
notes) are untested here -- it is precise evidence that \emph{this specific,
simple linear-in-mean family} is insufficient, reported as such rather than
rounded to a verdict either way, mirroring the family-by-family honesty
standard used throughout this research programme's other existence
certifications.

\subsection{Answering ``does the existing decoder admit the required \texorpdfstring{$\Q$}{Q}-dynamics''}

Not established by either family tested here (the existing trained
target-measure kernel, or a first-order additive mean-shift). Sec.~\ref{sec:sec5}
proceeds with the CRC sequence regardless: its own step (E), the constrained
inverse under exact price preservation, is a different, more tractable
object than a global dynamic no-arbitrage guarantee, and is the concrete
engineering gap the audit identified.

\section{The consistent-recalibration sequence, implemented}
\label{sec:sec5}

Implements steps A--G exactly as specified; the clearest gap identified in
Sec.~\ref{sec:audit} (no price-preserving reset existed anywhere in
\texttt{flex/} before this).

\subsection{A structural fact exploited, and verified}

\texttt{bound\_native\_raw} maps \emph{any} raw torch tensor to bounded,
admissible \texttt{(dlvsqs,forward\_vols,eta)}, and the SANOS forward
recursion then guarantees a non-negative discrete density from positive
\texttt{dlvsqs} -- verified directly: the minimum decoded density stayed at
$\approx0.0106$--$0.0107$ across every tested reset below, never approaching
zero. \textbf{Consequence}: for this decoder, step (E)'s ``constrained
inverse... subject to positivity... and other model admissibility
constraints'' is an \emph{unconstrained} optimisation in raw-state space --
admissibility is automatic by construction, not something the optimiser must
separately enforce.

\subsection{The sequence, run end to end}

4 resets, synthetic independent oracle standing in for the observed market
(see data-honesty note, Sec.~\ref{sec:deferred}), 15-contract protected book,
tolerance $10^{-4}$:

\begin{center}
\begin{tabular}{rlrrrr}
\toprule
Step & Converged & Iters & Pre-corr.\ $\max|\text{err}|$ & Post-corr.\ $\max|\text{err}|$ & Raw-state move ($L_2$) \\
\midrule
0 (initial calib.) & yes & 66 & 0.00272 & $8.8\times10^{-5}$ & 0.755 \\
1 & yes & 5 & 0.00021 & $7.2\times10^{-5}$ & 0.062 \\
2 & yes & 3 & 0.00038 & $2.0\times10^{-4}$ & 0.144 \\
3 & yes & 9 & 0.00021 & $1.0\times10^{-4}$ & 0.156 \\
\bottomrule
\end{tabular}
\end{center}

\textbf{Distinguishing the three things Sec.~5 asks to keep separate}:
\emph{stochastic market evolution} is the independent oracle state, stepped
under the fitted P law, never touched by the reset; \emph{price-preserving
representation change} is the reset representation's own P-law one-step
proposal (already tracking reasonably: pre-correction errors are
$20$--$30\times$ smaller than the initial uncalibrated guess) corrected by a
further optimisation an order of magnitude smaller in raw-state motion than
the initial calibration; \emph{conditioning on newly observed data} is
step~C's fresh evaluation of the oracle at its new state -- the reset never
assumes the market re-realises what the representation predicted.

\subsection{Feasibility decision is not vacuous}

A deliberately implausible target book (uniform random prices up to $3.0$,
tolerance $10^{-8}$, $20$ iterations) is correctly reported as
\textbf{not converged} (residual $2.82\to2.82$, essentially unmoved) --
confirming the feasibility check genuinely can, and does, reject, rather than
reporting ``accept'' trivially in every run.

\subsection{Scope of what this establishes}

\textbf{Does}: a complete, working, price-preserving reset sequence for the
FLEX decoder, with genuine feasibility accounting on both the success and the
failure side, and a clean separation of market evolution from representation
correction. \textbf{Does not}: establish that the resulting sequence of reset
representations is itself dynamically consistent in the Sec.~\ref{sec:sec23}
sense -- the constrained inverse guarantees only static admissibility and
exact price reproduction \emph{at} the reset instant, exactly the distinction
the assignment itself draws between agreement at a reset and coherence of the
continuation model. That question remains open per
Sec.~\ref{sec:sec23}/\ref{sec:deferred}.

\section{What is deferred, and why -- precisely scoped}
\label{sec:deferred}

\subsection{Data honesty, stated once more explicitly}

This environment's local \texttt{QUANTDATA\_TEST\_MODE=1} fixture data
(\texttt{tests/stored\_data/}) covers only a handful of scattered single
trading days per identifier (2005, 2010--2011, 2024, early 2025) -- not a
continuous multi-year weekly series. The original assignment's own train/
validation split (2022--2024 / 2025) requires the production
\texttt{QUANTDATA\_ROOT}/WRDS-backed history, which was not configured here.
Every number in Secs.~\ref{sec:sec23}--\ref{sec:sec5} therefore validates
\emph{machinery} (does the pricing identity hold, is the drift feasible, does
the reset converge and correctly detect infeasibility) on synthetic-but-
structurally-faithful data -- the same standard the repository's own unit
tests already use -- not real market performance.

\subsection{Not attempted here, and what would be required}

\begin{itemize}[leftmargin=1.4em]
\item \textbf{A demonstrably admissible $\Q$ family for FLEX.}
  Sec.~\ref{sec:sec23} rules out (for the tested instrument set) two
  specific families -- the existing trained target-measure kernel and a
  first-order additive mean-shift -- without ruling out richer ones (a full
  covariance tilt with a larger instrument set; an independent-product
  construction with a genuine conditional pricing kernel for a stochastic-
  volatility $Y$, per this research programme's other notes; a structural
  regridding of the smoothing/anchor coordinates themselves). Establishing
  or ruling out a richer family is the natural next step and was out of
  reach within this engagement's time budget once the audit, the pricing-
  identity formalisation, the feasibility ranking, and the CRC sequence were
  each done properly rather than superficially.
\item \textbf{Fitted, compatible $\Pp$ dynamics via risk premia}
  ($b_\Pp=b_\Q+\Sigma\lambda$, or a discrete positive-tilt kernel with valid
  density ratios). This presupposes an admissible $\Q$ family to tilt away
  from, which is not yet established (above); attempting it before that
  would risk exactly the error the assignment warns against -- reporting
  joint estimation as though it were CRC.
\item \textbf{Historical performance validation} (chronological train/
  validation split, quote fit, transition likelihood, covariance/
  autocorrelation/tail diagnostics, spot--surface dependence, OCE/hedging
  diagnostics) against real option data. Requires the production data
  environment.
\item \textbf{Reading the existing \texttt{run\_*\_q\_existence.py} result
  artefacts.} Their training-cache outputs are not present in this checkout.
  Flagged explicitly in the audit (Sec.~\ref{sec:audit}) rather than silently
  bypassed; reading them before extending the $\Q$-family search further is
  the responsible next step, per the assignment's own instruction not to
  silently duplicate or contradict unread prior work.
\end{itemize}

This is not reported as a universal impossibility result for any of the above
-- consistent with the assignment's own explicit instruction not to turn the
failure of one parametrisation, or the absence of one numerical search, into
such a claim. Each item above is a well-defined, scoped next step, not an
open-ended one.

\section{Methods and reproduction}
\label{sec:methods}

\textbf{Software}: Python 3.12 (\texttt{/opt/homebrew/bin/python3.12}), a
fresh virtual environment built via \texttt{python3.12 -m venv .venv312} at
the \texttt{quant\_data-main} repository root, dependencies installed via
\texttt{pip install -e .} against the repository's own
\texttt{pyproject.toml} (torch, cdxcore, and the rest of the declared stack
are all public PyPI packages; no private credentials were required).

\textbf{Reproduction}:
\begin{verbatim}
cd quant_data-main
python3.12 -m venv .venv312
.venv312/bin/python -m pip install -e . pytest pytest-cov

# baseline (untouched, must pass before trusting anything below)
QUANTDATA_TEST_MODE=1 .venv312/bin/python -m pytest \
    tests/test_dysanos_torch_simulator.py tests/test_dysanos_market.py -q

# this report's new code and tests
QUANTDATA_TEST_MODE=1 .venv312/bin/python -m pytest \
    tests/test_flex_crc_extension.py -q

QUANTDATA_TEST_MODE=1 .venv312/bin/python -m \
    quantdata.dysanos.flex.crc_extension.section2_pricing_identity
QUANTDATA_TEST_MODE=1 .venv312/bin/python -m \
    quantdata.dysanos.flex.crc_extension.section5_crc_sequence
\end{verbatim}

\textbf{New files} (nothing outside \texttt{flex/crc\_extension/} and
\texttt{tests/test\_flex\_crc\_extension.py} was modified):
\texttt{flex/crc\_extension/PHASE1\_AUDIT.md},
\texttt{SECTION2\_3\_FINDINGS.md}, \texttt{SECTION5\_FINDINGS.md},
\texttt{section2\_pricing\_identity.py}, \texttt{section5\_crc\_sequence.py},
this report.

\textbf{Seeds}: every reported number above used a declared, fixed seed
(\texttt{seed=17} for the primary martingale test, \texttt{seed=314} for
target-kernel training, \texttt{seed=71} for the feasibility Jacobian,
\texttt{seed=2024} for the CRC sequence) -- re-running the commands above
reproduces the same numbers bit-for-bit (torch's CPU RNG under
\texttt{manual\_seed}).

\textbf{Reused, unmodified modules}: \texttt{weekly\_dynamics.py}
(\texttt{fit\_weekly\_hierarchical\_dynamics}, \texttt{TorchWeeklyDynamics}),
\texttt{weekly\_simulation.py} (\texttt{WeeklySanosSimulator},
\texttt{hjm\_roll\_residuals}, \texttt{hjm\_roll\_moment\_loss}),
\texttt{weekly\_measure.py} (\texttt{NormalWeeklyMeasureKernel}),
\texttt{model.py} (\texttt{bound\_native\_raw}, \texttt{LatentStateSanosModel}),
\texttt{sanos.py} (\texttt{FlexSanosDecoder}, \texttt{price\_calls\_batch}).


\section{Conclusion}
\label{sec:conclusion}

The FLEX simulator's existing statistical (P) dynamics do not, and were never
claimed to, satisfy a risk-neutral pricing identity; this is now verified with a
well-powered numerical test rather than left as an unverified docstring
assertion. Its existing gradient-trained target-measure machinery narrows the
gap but does not close it for the tested family, and a first-order feasibility
analysis shows why: a simple additive drift adjustment is both rank-deficient
and insufficient even on its own reachable subspace. Independently of that open
question, the concrete engineering gap identified by the audit -- an exact,
feasibility-checked, price-preserving reset -- is closed: the CRC sequence
implemented here (Sec.~\ref{sec:sec5}) demonstrably calibrates, evolves,
proposes, inverts, decides feasibility, and continues, exactly as specified,
with a genuine (not vacuous) reject path verified. Closing the remaining gap --
a demonstrably admissible $\Q$ family, fitted historical $\Pp$ dynamics compatible
with it, and validation against real option data -- is scoped precisely in
Sec.~\ref{sec:deferred} as the next, well-defined step, not asserted here and not
claimed impossible either.

\end{document}
